\subsection{\texorpdfstring{The homomorphism \(\operatorname{Ext}_{\mathbf{A}}(\mathbf{P}, \mathbf{Q}) \otimes \operatorname{Tor}^{\mathbf{A}}(\mathbf{P}, \mathbf{M}) \to \operatorname{Tor}^{\mathbf{A}}(\mathbf{Q}, \mathbf{M})\)}{The homomorphism \textbackslash operatorname\{Ext\}\_\{\textbackslash mathbf\{A\}\}(\textbackslash mathbf\{P\}, \textbackslash mathbf\{Q\}) \textbackslash otimes \textbackslash operatorname\{Tor\}\^{}\{\textbackslash mathbf\{A\}\}(\textbackslash mathbf\{P\}, \textbackslash mathbf\{M\}) \textbackslash to \textbackslash operatorname\{Tor\}\^{}\{\textbackslash mathbf\{A\}\}(\textbackslash mathbf\{Q\}, \textbackslash mathbf\{M\})}}

Let M be a left A-module, and let P and Q be two right A-modules. Consider the homomorphism \(\operatorname{Homgr}_{\mathbf{A}}(\mathrm{L}(\mathrm{P}), \mathrm{L}(\mathrm{Q})) \otimes_{k} (\mathrm{L}(\mathrm{P}) \otimes_{\mathbf{A}} \mathrm{L}(\mathrm{M})) \to \mathrm{L}(\mathrm{Q}) \otimes_{\mathbf{A}} \mathrm{L}(\mathrm{M})\) which to \(f \otimes (x \otimes y)\) associates \(f(x) \otimes y\) . By X, p.~99, this is a morphism of complexes. From this one deduces a graded \(k\) -linear map of degree 0

\[
\mathrm{H} (\operatorname{Homgr} _ {\mathrm{A}} (\mathrm{L} (\mathrm{P}), \mathrm{L} (\mathrm{Q}))) \otimes_ {k} \operatorname{Tor} ^ {\mathrm{A}} (\mathrm{P}, \mathrm{M}) \rightarrow \operatorname{Tor} ^ {\mathrm{A}} (\mathrm{Q}, \mathrm{M})
\]

hence, via the isomorphism \(\varphi(\mathbf{L}(\mathbf{P}),\mathbf{L}(\mathbf{Q}))\) of § 6 (X, p.~100, th. 1), a graded \(k\) -linear map of degree 0

\[
\operatorname{Ext} _ {\mathbf {A}} (\mathrm{P}, \mathrm{Q}) \otimes_ {k} \operatorname{Tor} ^ {\mathbf {A}} (\mathrm{P}, \mathrm{M}) \rightarrow \operatorname{Tor} ^ {\mathbf {A}} (\mathrm{Q}, \mathrm{M}),\tag{10}
\]

corresponding to k-bilinear maps

\[
c _ {\mathrm{P}, \mathrm{Q}; \mathrm{M}}: \operatorname{Ext} _ {\mathrm{A}} ^ {n} (\mathrm{P}, \mathrm{Q}) \times \operatorname{Tor} _ {m} ^ {\mathrm{A}} (\mathrm{P}, \mathrm{M}) \rightarrow \operatorname{Tor} _ {m - n} ^ {\mathrm{A}} (\mathrm{Q}, \mathrm{M});\tag{11}
\]

the image of the pair \((\alpha, \gamma)\) by \(c_{P,Q;M}\) is called the composition product of \(\alpha\) and \(\gamma\) and is denoted \(\alpha \circ \gamma\) .

By construction, \(\alpha \circ \gamma\) is obtained as follows: one represents \(\alpha\) by a morphism of complexes \(f: \mathrm{L}(\mathrm{P}) \to \mathrm{L}(\mathrm{Q})(-n)\) , \(\gamma\) by an element \(z \in Z_m(\mathrm{L}(\mathrm{P}) \otimes_{\mathbf{A}} \mathrm{L}(\mathrm{M}))\) , and \(\alpha \otimes \gamma\) is the class of the element

\[
(f \otimes 1) (z) \in Z _ {m} (\mathrm{L} (\mathrm{Q}) (- n) \otimes_ {\mathrm{A}} \mathrm{L} (\mathrm{M})) = Z _ {m - n} (\mathrm{L} (\mathrm{Q}) \otimes_ {\mathrm{A}} \mathrm{L} (\mathrm{M})).
\]

For example, if \(\alpha \in \operatorname{Hom}_{\mathbf{A}}(\mathbf{P},\mathbf{Q})\) , then \(\alpha \circ \gamma = \operatorname{Tor}(\alpha,1)(\gamma)\) .

Remarks. --- 1) If one uses the isomorphisms \(\psi\) of X, p.~69, one can also define the composition product by the commutative diagram

\[
\begin{array}{c} \operatorname{Ext} _ {\mathrm{A}} ^ {n} (\mathrm{P}, \mathrm{Q}) \times \operatorname{Tor} _ {m} ^ {\mathrm{A}} (\mathrm{P}, \mathrm{M}) \xrightarrow {c _ {\mathrm{P} , \mathrm{Q} : \mathrm{M}}} \operatorname{Tor} _ {m - n} ^ {\mathrm{A}} (\mathrm{Q}, \mathrm{M}) \\ \bar {a} _ {\mathrm{P}, \mathrm{Q}} \times \psi_ {\mathrm{P}} (\mathrm{M}) \Biggl \downarrow \\ \mathrm{H} ^ {n} (\operatorname{Homgr} _ {\mathrm{A}} (\mathrm{L} (\mathrm{P}), \mathrm{L} (\mathrm{Q}))) \times \mathrm{H} _ {m} (\mathrm{L} (\mathrm{P}) \otimes_ {\mathrm{A}} \mathrm{M}) \longrightarrow \mathrm{H} _ {m - n} (\mathrm{L} (\mathrm{Q}) \otimes_ {\mathrm{A}} \mathrm{M}); \end{array}
\]

in other words, we represent \(\alpha\) by a morphism \(f\) from L(P) to L(Q) \((-n)\) , \(\gamma\) by a cycle \(x \in \mathrm{L}_{m}(\mathrm{P}) \otimes_{\mathbf{A}} \mathrm{M}\) , and \(\alpha \circ \gamma\) is the cycle class of

\[
\left(f _ {m} \otimes 1 _ {\mathbf {M}}\right) (x) \in \mathrm{L} _ {m - n} (\mathrm{Q}) \otimes_ {\mathbf {A}} \mathrm{M}.
\]

\begin{enumerate}
\def\labelenumi{\arabic{enumi})}
\setcounter{enumi}{1}
\tightlist
\item
  One can also use the resolutions I(P) and I(Q).
\end{enumerate}

Similarly, if N is a second left A-module, one defines a composition product \((\mu, \gamma) \mapsto \mu \circ \gamma\) , denoted

\[
c _ {\mathrm{P}; \mathrm{M}, \mathrm{N}}: \operatorname{Ext} _ {\mathrm{A}} ^ {r} (\mathrm{M}, \mathrm{N}) \times \operatorname{Tor} _ {m} ^ {\mathrm{A}} (\mathrm{P}, \mathrm{M}) \rightarrow \operatorname{Tor} _ {m - r} ^ {\mathrm{A}} (\mathrm{P}, \mathrm{N})
\]

by the commutative diagram

\[
\begin{array}{c} \operatorname{Ext} _ {\mathbf {A}} ^ {r} (\mathbf {M}, \mathbf {N}) \times \operatorname{Tor} _ {m} ^ {\mathbf {A}} (\mathbf {P}, \mathbf {M}) \xrightarrow {c _ {\mathrm{P} ; \mathrm{M} , \mathrm{N}}} \operatorname{Tor} _ {m - r} ^ {\mathbf {A}} (\mathbf {P}, \mathbf {N}) \\ 1 \times \sigma_ {\mathrm{P}, \mathrm{M}, r} \Biggl \downarrow \\ \operatorname{Ext} _ {\mathbf {A} ^ {\circ}} ^ {r} (\mathbf {M} ^ {\circ}, \mathbf {N} ^ {\circ}) \times \operatorname{Tor} _ {m} ^ {\mathbf {A} ^ {\circ}} (\mathbf {M} ^ {\circ}, \mathbf {P} ^ {\circ}) \xrightarrow {c _ {\mathrm{M} ^ {\circ} , \mathrm{N} ^ {\circ} ; \mathrm{P} ^ {\circ}}} \operatorname{Tor} _ {m - r} ^ {\mathbf {A} ^ {\circ}} (\mathbf {N} ^ {\circ}, \mathbf {P} ^ {\circ}) \end{array}\tag{12}
\]

where σ denotes the commutation isomorphisms (X, p.~71).

If \(\mu \in \operatorname{Ext}_{\mathbf{A}}^{r}(\mathbf{M}, \mathbf{N})\) is the class of the morphism \(g: \mathrm{L}(\mathbf{M}) \to \mathrm{L}(\mathbf{N})(-r)\) , and if \(\gamma \in \operatorname{Tor}_{m}^{\mathbf{A}}(\mathbf{P}, \mathbf{M})\) is the cycle class of \(z = \sum_{i,j} z_{ij}\) , where \(z_{ij} \in \mathrm{L}_i(\mathbf{P}) \otimes_{\mathbf{A}} \mathrm{L}_j(\mathbf{M})\) , \(\mu \circ \gamma\) is therefore the class of the cycle \(\sum_{i,j} (-1)^{ir}(1 \otimes g)(z_{ij})\) .

One can also represent \(\gamma\) by a cycle \(y \in \mathrm{P} \otimes \mathrm{L}_{m}(\mathbf{M})\) , and \(\mu \circ \gamma\) is the cycle class of \((1 \otimes g)(y) \in \mathrm{P} \otimes \mathrm{L}_{m-r}(\mathbf{M})\) .

PROPOSITION 6. --- Let K, M, N be left A-modules, and P, Q, R right A-modules,

\[
\alpha \in \operatorname{Ext} _ {\mathrm{A}} ^ {n} (\mathrm{P}, \mathrm{Q}), \beta \in \operatorname{Ext} _ {\mathrm{A}} ^ {p} (\mathrm{Q}, \mathrm{R}), \lambda \in \operatorname{Ext} _ {\mathrm{A}} ^ {r} (\mathrm{K}, \mathrm{M}), \mu \in \operatorname{Ext} _ {\mathrm{A}} ^ {s} (\mathrm{M}, \mathrm{N}), \gamma \in \operatorname{Tor} _ {m} ^ {\mathrm{A}} (\mathrm{P}, \mathrm{K}).
\]

Then

\begin{enumerate}
\def\labelenumi{(\arabic{enumi})}
\setcounter{enumi}{12}
\tightlist
\item
\end{enumerate}

\[
(\beta \circ \alpha) \circ \gamma = \beta \circ (\alpha \circ \gamma) \text { in }\operatorname{Tor} _ {m - p - n} ^ {\mathrm{A}} (\mathrm{R}, \mathrm{K}),\tag{14}
\]

\[
(\mu \circ \lambda) \circ \gamma = \mu \circ (\lambda \circ \gamma) \text { in }\operatorname{Tor} _ {n - r - s} ^ {\mathbf {A}} (\mathrm{P}, \mathrm{N}),\tag{15}
\]

\[
\alpha \circ (\lambda \circ \gamma) = (- 1) ^ {n r} \lambda \circ (\alpha \circ \gamma) \text { in }\operatorname{Tor} _ {m - p - r} ^ {\mathrm{A}} (\mathrm{Q}, \mathrm{M}).
\]

Formulas (13) and (14) follow immediately from the definitions. Let us prove (15). Let \(z = \sum z_{ij}, z_{ij} \in \mathbf{L}_i(\mathbf{P}) \otimes \mathbf{L}_j(\mathbf{K})\) a cycle representing \(\gamma, f: \mathrm{L}(\mathrm{P}) \to \mathrm{L}(\mathrm{Q})(-n)\) and \(g: \mathbf{L}(\mathbf{K}) \to \mathbf{L}(\mathbf{M})(-r)\) morphisms representing \(\alpha\) and \(\lambda\) . Then \(\lambda \circ (\alpha \circ \gamma)\) is the class of \(\sum (-1)^{(i-n)r}(f \otimes g)(z_{ij})\) and \(\alpha \circ (\lambda \circ \gamma)\) is the class of

\[
\sum (- 1) ^ {i r} (f \otimes g) (z _ {i j}), \quad \text { whence } (1 5).
\]

